\newpage
\begin{center}
\large\textbf{Минобрнауки России}

\large\textbf{Юго-Западный государственный университет}
\vskip 1em
\normalsize{Кафедра программной инженерии}
\vskip 1em
\ifВКР{
        \begin{flushright}
        \begin{tabular}{p{.4\textwidth}}
        \centrow УТВЕРЖДАЮ: \\
        \centrow Заведующий кафедрой \\
        \hrulefill \\
        \setarstrut{\footnotesize}
        \centrow\footnotesize{(подпись, инициалы, фамилия)}\\
        \restorearstrut
        «\underline{\hspace{1cm}}»
        \underline{\hspace{3cm}}
        20\underline{\hspace{1cm}} г.\\
        \end{tabular}
        \end{flushright}
        }\fi
\end{center}
\vspace{1em}
  \begin{center}
  \large
\ifВКР{
ЗАДАНИЕ НА ВЫПУСКНУЮ КВАЛИФИКАЦИОННУЮ РАБОТУ
  ПО ПРОГРАММЕ БАКАЛАВРИАТА}
  \else
ЗАДАНИЕ НА КУРСОВУЮ РАБОТУ (ПРОЕКТ)
\fi
\normalsize
  \end{center}
\vspace{1em}
{\parindent0pt
  Студента \АвторРод, шифр\ \Шифр, группа \Группа
  
1. Тема «\Тема\ \ТемаВтораяСтрока»
\ifВКР{
утверждена приказом ректора ЮЗГУ от \ДатаПриказа\ № \НомерПриказа
}\fi.

2. Срок предоставления работы к защите \СрокПредоставления

3. Исходные данные для создания программной системы:

3.1. Перечень решаемых задач:}

\renewcommand\labelenumi{\theenumi)}

\begin{enumerate}
\item Провести анализ бизнес-процессов компании в области управления персоналом и внутренних коммуникаций.
\item Разработать концептуальную модель веб-платформы, объединяющей почту, календарь, чат, управление задачами и видеоконференции.
\item Спроектировать архитектуру клиент-серверного приложения на основе современных веб-технологий (TypeScript, React, Node.js, JMAP).
\item Разработать интерфейсные модули для основных сервисов: «Почта», «Контакты», «Проекты», «Календарь», «Файлы», «Разговоры», «Видеоконференции», «Настройки».
\item Реализовать серверную часть платформы и интеграцию с внешними сервисами (почта, хранилище, база данных).
\item Провести модульное и интеграционное тестирование платформы.
\item Внедрить веб-платформу в инфраструктуру организации и оценить её эффективность.
\end{enumerate}

{\parindent0pt
  3.2. Входные данные и требуемые результаты для программы:}

\begin{enumerate}
\item Входными данными для программной системы являются: информация о сотрудниках; структура задач и проектов; события календаря; входящие и исходящие сообщения электронной почты; файлы и документы; данные видеоконференций; пользовательские настройки интерфейса и уведомлений.
\item Выходными данными для программной системы являются: расписание задач и мероприятий, отображаемое в интерфейсе панели управления; статусы задач, комментарии, вложения и история выполнения; электронные письма и вложения с возможностью фильтрации и поиска; отчёты по активности сотрудников, прогрессу выполнения задач, загрузке хранилища; сводные панели и виджеты по текущей работе, уведомления и оповещения; сведения о прошедших видеовстречах и сохранённых материалах; настроенные профили пользователей, включая тему, язык, подпись.
\end{enumerate}

{\parindent0pt

  4. Содержание работы (по разделам):
  
  4.1. Введение.

  4.2. Анализ предметной области.

4.3. Техническое задание: цель, назначение, функциональные и нефункциональные 
требования, требования к интерфейсу и документации.

4.4. Технический проект: концептуальная модель системы, архитектура веб-платформы,
проектирование интерфейсов и базы данных.

4.5. Рабочий проект: реализация клиентской и серверной частей, описание ключевых 
модулей, тестирование (модульное, интеграционное), интеграция с внешними сервисами.

4.6. Заключение

4.7. Список использованных источников

5. Перечень графического материала:

\списокПлакатов

\vskip 2em
\begin{tabular}{p{6.8cm}C{3.8cm}C{4.8cm}}
Руководитель \ifВКР{ВКР}\else работы (проекта) \fi & \lhrulefill{\fill} & \fillcenter\Руководитель\\
\setarstrut{\footnotesize}
& \footnotesize{(подпись, дата)} & \footnotesize{(инициалы, фамилия)}\\
\restorearstrut
Задание принял к исполнению & \lhrulefill{\fill} & \fillcenter\Автор\\
\setarstrut{\footnotesize}
& \footnotesize{(подпись, дата)} & \footnotesize{(инициалы, фамилия)}\\
\restorearstrut
\end{tabular}
}

\renewcommand\labelenumi{\theenumi.}
